\par\noindent\begin{minipage}{\linewidth}
\subsection{Grenzwerte im Unendlichen}
\label{subsec:real_functions_infinite_limits}
Häufig interessiert uns das asymptotische Verhalten einer Funktion,
also wie sie sich für betragsmäßig große $x$ verhält.
Dieses Verhalten beschreiben wir mit Grenzwerten für $x\to\pm\infty$.
\begin{DefinitionBox}{Grenzverhalten im Unendlichen}{real_functions_limit_at_infinity}
  Wächst $x$ unbegrenzt ins Positive (oder fällt unbegrenzt ins Negative),
  also $x\to+\infty$ (bzw. $x\to-\infty$), und nähern sich dabei die
  Funktionswerte $f(x)$ einem Wert $L\in\R$, so heißt $L$ der
  \emph{Grenzwert von $f$ im Unendlichen}. Wir schreiben:
  \[
    \lim_{x\to+\infty}f(x)=L
    \qquad\left(\text{bzw. }\lim_{x\to-\infty}f(x)=L\right).
  \]
\end{DefinitionBox}
\begin{ExampleBox}[before skip=0.3\baselineskip]{Grenzverhalten einer Geraden}{real_functions_identity_infinity}
  Am Graphen der Geraden betrachten wir das Verhalten für $x\to\pm\infty$.
  \[
    \lim_{x\to\pm\infty}f(x).
  \]
  \par\smallskip\noindent
  \begin{minipage}[c]{0.47\linewidth}
    Für $x\to+\infty$ wachsen die Funktionswerte unbegrenzt ins Positive;
    für $x\to-\infty$ fallen sie unbegrenzt ins Negative.
    \[
      \begin{aligned}
        \lim_{x\to+\infty}x&=\Cloze{1.6cm}{+\infty},\\[0.5em]
        \lim_{x\to-\infty}x&=\Cloze{1.6cm}{-\infty}.
      \end{aligned}
    \]
    Beide Grenzwerte sind uneigentlich.
  \end{minipage}\hfill
  \begin{minipage}[c]{0.47\linewidth}
    \centering
    \CachedTikzFigure{content/resources/chapters/real_functions/tikz/functions_identity_infinity.tex}
  \end{minipage}
\end{ExampleBox}
\end{minipage}
